\documentclass[aspectratio=169,10pt,english]{beamer}

\input{../../../_preambulo_beamer.tex}
\input{../../../_comandos_beamer.tex}

\selectlanguage{english}

\title{MA1001B - Statistical Modeling for Decision Making}
\subtitle{Lesson 34: Z Test for a Mean with Known Variance}
\author{}
\institute{Tecnol\'ogico de Monterrey}
\date{}

\begin{document}

\begin{frame}[plain]
  \vspace{1.2cm}
  {\Large\bfseries MA1001B - Statistical Modeling for Decision Making\par}
  \vspace{0.7cm}
  {\large Lesson 34: Z Test for a Mean with Known Variance\par}
  \vspace{0.7cm}
  {\color{TecRojo}\rule{\textwidth}{1.2pt}}
  \vspace{0.8cm}

  MA1001B Faculty\\[0.3cm]
  Term\\[0.3cm]
  Tecnol\'ogico de Monterrey
\end{frame}

\begin{frame}{The Known-Sigma Test Question}
  \begin{alertblock}{Guiding question}
    How is a one-sided test of an unknown mean carried out when $\sigma$ is
    treated as known?
  \end{alertblock}

  \vspace{0.6em}
  By the end of this lesson, you can:
  \begin{itemize}
    \item compute $z=(\bar{x}-\mu_{0})/(\sigma/\sqrt{n})$;
    \item obtain the upper-tail $p$-value from the standard normal;
    \item decide at $\alpha=0.05$ by comparing $p$ with $\alpha$;
    \item explain why known $\sigma$ is rarely true in operations data.
  \end{itemize}

  \vfill
  \scriptsize All delivery times used today are synthetic. No real company data are used.
\end{frame}

\begin{frame}{A Synthetic Delivery-Time Test}
  \small
  \begin{center}
    \begin{tabular}{lr}
      \toprule
      \textbf{Quantity} & \textbf{Value} \\
      \midrule
      $H_{0}$ & $\mu=80$ minutes \\
      $H_{1}$ & $\mu>80$ minutes \\
      $n$, known $\sigma$ & 40, 6 minutes \\
      Sample $\bar{x}$ (seed 42) & 82.229590 minutes \\
      \bottomrule
    \end{tabular}
  \end{center}

  \vspace{0.5em}
  \begin{exampleblock}{Standard error}
    \[
      \mathrm{SE}=\frac{6}{\sqrt{40}}=0.948683.
    \]
    The analyst treats $\sigma$ as known from long-run process capability.
  \end{exampleblock}
\end{frame}

\begin{frame}[fragile]{Step 1: The Z Statistic}
  \begin{columns}[T]
    \begin{column}{0.42\textwidth}
      \small
      \begin{lstlisting}
se = sigma / np.sqrt(n)
z = (xbar - 80) / se
      \end{lstlisting}

      \begin{block}{Verified output}
        $\bar{x}=82.229590$\\
        $\mathrm{SE}=0.948683$\\
        $z=2.350194$
      \end{block}
    \end{column}
    \begin{column}{0.58\textwidth}
      \centering
      \includegraphics[width=\textwidth]{../figures/z_test_mean_01_z_statistic.png}
      \vspace{0.2em}
      \scriptsize Delivery histogram from Step 1
    \end{column}
  \end{columns}
\end{frame}

\begin{frame}{Upper-Tail $P$-Value and Decision}
  \small
  \[
    p=P(Z\ge 2.350194)=\texttt{norm.sf}(2.350194)=0.009382.
  \]
  Because $p<\alpha=0.05$, $H_{0}$ is rejected: the synthetic sample is
  inconsistent with a mean of 80 minutes when $\sigma=6$ is treated as known.

  \vspace{0.4em}
  The $p$-value is not the probability that $H_{0}$ is true. It is the tail
  probability of a result this large or larger under $H_{0}$.
\end{frame}

\begin{frame}[fragile]{Step 2: $P$-Value and Decision}
  \begin{columns}[T]
    \begin{column}{0.40\textwidth}
      \small
      \begin{lstlisting}
p = stats.norm.sf(z)
reject = p < 0.05
      \end{lstlisting}

      \begin{block}{Verified output}
        $p=0.009382$\\
        $\alpha=0.05$\\
        Reject $H_{0}$
      \end{block}
    \end{column}
    \begin{column}{0.60\textwidth}
      \centering
      \includegraphics[width=\textwidth]{../figures/z_test_mean_02_p_value.png}
      \vspace{0.2em}
      \scriptsize Standard-normal tail from Step 2
    \end{column}
  \end{columns}
\end{frame}

\begin{frame}{Step 3: Known $\sigma$ Is Rarely True}
  \begin{columns}[T]
    \begin{column}{0.42\textwidth}
      \small
      Sample $s=4.955099\neq 6$.

      \vspace{0.4em}
      Invalid z-with-$s$:
      \[
        z=2.845789,\quad p=0.002215.
      \]

      \vspace{0.5em}
      \textbf{Limit:} keeping a $z$ curve after replacing $\sigma$ by $s$ is
      the wrong test. Lesson 35 uses $t$ with $\mathrm{df}=n-1$.
    \end{column}
    \begin{column}{0.58\textwidth}
      \centering
      \includegraphics[width=\textwidth]{../figures/z_test_mean_03_known_sigma_limit.png}
      \vspace{0.2em}
      \scriptsize Valid z versus invalid z-with-$s$ from Step 3
    \end{column}
  \end{columns}
\end{frame}

\begin{frame}{Concept Check}
  \begin{enumerate}
    \item Compute $\mathrm{SE}=6/\sqrt{40}$.
    \item Compute $z=(82.229590-80)/\mathrm{SE}$.
    \item Why is $p=0.009382$ not the probability that $\mu=80$?
    \item Why is $s=4.955099$ not a license to keep the z test?
  \end{enumerate}

  \vspace{0.6em}
  \begin{block}{Connection to Lesson 35}
    The session can continue directly into the one-sample $t$ test, with $s$
    in the denominator and $\mathrm{df}=n-1$.
  \end{block}

  \vfill
  \scriptsize Source: OpenStax, \textit{Introductory Business Statistics 2e},
  Chapter 9, Sections 9.3--9.4. CC BY-NC-SA.
\end{frame}

\appendix



\end{document}
